\documentclass[11pt, a4paper]{../problems_template}

\begin{document}

\begin{titlepage}
  \begin{center}
	\Huge{Конкурс имени Савина в журнале Квант}
  \end{center}
\end{titlepage} 

\chapter{2017-2018}

\begin{exercise}[7]
Все натуральные числа раскрасили в сто цветов. Докажите, что найдется несколько (не менее двух) различных чисел одного цвета такие, что их произведение имеет ровно 1000 различных натуральных делителей.
\end{exercise}

\begin{exercise}
Кузнечик умеет прыгать по полоске из $n$ клеток на 8, 9 и 10 клеток в любую сторону. Будем называть натуральное число $n$ пропрыгиваемым, если кузнечик может, начав с некоторой клетки, обойти всю полоску, побывав на каждой клетке ровно один раз. Докажите, что начиная с некоторого числа $M$, все $n > M$ пропрыгиваемы.
\end{exercise}

\chapter{2018-2019}

\begin{exercise}[3]
Квантик и Ноутик показывают такой фокус. Зритель задумывает любые шесть разных целых чисел от 1 до 125 и сообщает их только Ноутику. После этого Ноутик называет Квантику какие-то пять из них, и Квантик угадывает шестое задуманное зрителем число. Предложите способ, как могли бы действовать Квантик и Ноутик, чтобы фокус всегда удавался.
\end{exercise}

\begin{exercise}[4]
Квадрат $100 \times 100$ разбит на доминошки. Можно ли гарантировать, что его можно разрезать на две связные части, симметричные друг другу относительно центра квадрата, не разрезав ни одну доминошку? (Часть называется связной, если от любой ее клетки можно дойти до любой другой, совершая переходы в соседнюю по стороне клетку и не выходя за границу фигуры.)
\end{exercise}

\begin{exercise}
В клетки прямоугольника $2 \times n$ записывают числа $1, 2, \dots, 2n$ (в каждую клетку - по одному числу). Сначала в левую нижнюю угловую клетку записывают число 1, затем в одну из соседних
клеток - число 2, затем в одну из соседних с ранее занятыми - 3 и так далее все последующие числа записывают в клетки, соседние с ранее занятыми. Докажите, что существует ровно $2^{n - 1} n!$ способов
заполнить таблицу.
\end{exercise}

\chapter{2019-2020}

\begin{exercise}[16]
Квантик и Ноутик хотят показать такой фокус. Зритель задумывает два натуральных числа: $n$ и $n^{2}$, и сообщает одно из них Квантику, а другое - Ноутику. После этого Квантик показывает Ноутику чёрную или белую карточку, и Ноутик сразу угадывает число Квантика. Помогите Квантику и Ноутику договориться о своих действиях, чтобы фокус всегда удавался.
\end{exercise}

\chapter{2020-2021}

\begin{exercise}
Можно ли все натуральные числа окрасить в семь цветов так, чтобы каждый цвет присутствовал и произведение любых двух чисел одного цвета (в том числе, равных чисел) было числом того же цвета?
\end{exercise}

\begin{exercise}
Назовем последовательность из 50 натуральных чисел с суммой 100 хорошей, если из нее нельзя выбрать несколько подряд идущих чисел с суммой 50. Найдите количество таких последовательностей.
\end{exercise}

\chapter{2021-2022}

\begin{exercise}
\begin{enumerate}
\item Можно ли разрезать какой-нибудь прямоугольник на несколько равнобедренных прямоугольных \\ треугольников, среди которых нет одинаковых?
\item Можно ли так разрезать квадрат?
\end{enumerate}
\end{exercise}

\begin{exercise}[8]
Натуральное число называется палиндромом, если оно читается слева направо и справо налево одинаково \\ (например, 2, 33 или 12321). Для каких натуральных $n$ существует палиндром, делящийся на $n$?
\end{exercise}

\begin{exercise}[19]
Правильный треугольник со стороной 3 и правильный треугольник со стороной 4 в пересечении дают выпуклый шестиугольник периметра 7. Докажите, что у треугольников соответствующие стороны параллельны.
\end{exercise}

\chapter{2022-2023}

\begin{exercise}[12]
Числа $1, 2, \dots, 100$ покрасили в 10 цветов, в каждый цвет по 10 чисел. Затем посчитали все положительные разности между числами разного цвета. Какое наименьшее значение может принимать сумма всех таких \\ разностей?
\end{exercise}

\begin{exercise}[30]
Фокусник хочет заготовить 10 карточек, написать на каждой натуральное число, не большее 90, чтобы все числа были различны, и показывать такой фокус: зритель наугад выбирает две карточки,
называет фокуснику сумму чисел на них, а фокусник тут же отгадывает, какие две карточки у зрителя. Помогите фокуснику найти числа и объясните, почему фокус будет получаться.
\end{exercise}

\chapter{2023-2024}

\begin{exercise}[6]
Имеются стакан кофе, наполненный на $\frac{2}{3}$, и такой же стакан молока, наполненный на $\frac{2}{3}$. Разрешается переливать любое количество жидкости туда и обратно, тщательно её перемешивая, но нельзя ничего выливать. Можно ли получить в одном из стаканов напиток, составленный из молока и кофе в пропорции $1 : 1$?
\end{exercise}

\chapter{2024-2025}

\begin{exercise}[2]
У Пети есть набор из трех белых гирек массами 101 г, 102 г и 103 г, и такой же набор из трех черных гирек. Массы на гирьках не написаны, а на вид нельзя понять, какая гирька какой тяжелее. Петя хочет разбить гирьки на пары одинаковых по массе. Как ему сделать это за два взвешивания на чашечных весах со стрелкой, показывающих, какая чаша перевесила и на сколько грамм?
\end{exercise}

\begin{exercise}[8]
Клеточный прямоугольник по линиям сетки разрезан на полоски – прямоугольники шириной 1 и длиной больше 1. В каждой полоске две крайние клетки покрашены в красный цвет. Докажите, что от каждого угла \\ прямоугольника есть путь по красным клеткам в какой-то другой угол (путь можно пройти, переходя из клетки в соседнюю с ней по стороне).
\end{exercise}

\begin{exercise}[18]
Набор домино состоит из 28 различных прямоугольников $1 \times 2$, в клетках которых поставлено от 0 до 6 точек. Петя сложил все доминошки в произвольном порядке в кольцо так, что получилась прямоугольная рамка толщиной в клетку доминошки. Затем Вася склеил все доминошки по соседним сторонам, а потом разрезал каждую доминошку на две половинки. Могло ли оказаться, что полученные Васей доминошки тоже образуют полный набор?
\end{exercise}
 
\end{document}